\documentclass[11pt,a4paper]{article}
\usepackage[T1]{fontenc}
\usepackage[utf8]{inputenc}
\usepackage{lmodern}
\usepackage{amsmath,amssymb,amsthm}
\usepackage{booktabs}
\usepackage{xcolor}
\usepackage[margin=27mm]{geometry}
\usepackage[colorlinks=true,urlcolor=blue,citecolor=blue,linkcolor=blue]{hyperref}
\newtheorem{theorem}{Theorem}
\newtheorem{proposition}[theorem]{Proposition}
\newtheorem{lemma}[theorem]{Lemma}
\newtheorem{corollary}[theorem]{Corollary}
\theoremstyle{definition}
\newtheorem{definition}[theorem]{Definition}
\newtheorem{remark}[theorem]{Remark}
\theoremstyle{plain}
\DeclareMathOperator{\Tr}{Tr}
\newcommand{\ket}[1]{\lvert#1\rangle}
\newcommand{\bra}[1]{\langle#1\rvert}
\newcommand{\ketbra}[2]{\lvert#1\rangle\langle#2\rvert}
\setlength{\emergencystretch}{2em}
\setlength{\parindent}{0pt}
\setlength{\parskip}{0.3em plus 0.1em minus 0.1em}
\title{Separable states with the low-order moments of\\
states of arbitrary Schmidt rank at bounded energy}
\author{Fable \& Astra\\[0.35em]
\normalsize Research led by Jean-Philippe Bruneton\\[-0.1em]
\small Universit\'e Paris Cit\'e, France}
\date{19 November 2026}

\begin{document}
\maketitle
\begin{abstract}
In continuous-variable experiments the entanglement of two light modes is
usually assessed from the first and second moments of the field quadratures,
the mean vector and the covariance matrix, and recent work derives lower
bounds on the Schmidt number, the number of entangled degrees of freedom, from
these moments alone. We ask how much such data can say about the Schmidt
number when the mean photon number is bounded, and show that they cannot bound
it uniformly. For every $r\ge2$ we give a pure state of Schmidt rank $r$ and a
separable state with exactly the same first and second moments, and indeed the
same third moments, whose common mean photon number is at most a prescribed
bound independent of $r$. The two states are supported on the Fock vectors
$\ket{2j,2j}$, $0\le j<r$, the first as a coherent superposition and the second
as the corresponding classical mixture. Replacing the spacing $2$ by an
integer $m$ makes all quadrature moments of degree below $2m$ coincide, and a
moment of degree $2m$ distinguishes the states, so the construction is sharp.
At bounded energy the obstruction forces the $r$-th largest Schmidt coefficient
to be small, which points to a better-posed question.
\end{abstract}

\section{Introduction}
Entanglement between two modes of light is most often certified from
quadrature statistics. For Gaussian states the first and second moments
determine the state, and criteria built on the covariance matrix detect
entanglement and quantify it~\cite{Duan2000,Simon2000}. A finer question
concerns the Schmidt number~\cite{Terhal2000}, the smallest $k$ such that a
state is a mixture of pure states of Schmidt rank at most $k$, which counts how many degrees of freedom are genuinely entangled. Mishra,
Morelli, Baksov\'a, Rembold and Friis~\cite{Mishra2026} recently derived lower
bounds on the Schmidt number of continuous-variable states from their first and
second moments. It is natural to ask how far these moments can go: can a
criterion based on the covariance matrix, perhaps complemented by a bound on
the mean photon number, certify an arbitrarily large Schmidt number in a way
that applies uniformly to all states?

The simplest example shows that it cannot. Take the two-mode states
\[
\ket{\psi_r}=\sum_{j=0}^{r-1}\sqrt{p_j}\,\ket{2j,2j},
\qquad
\sigma_r=\sum_{j=0}^{r-1}p_j\,\ketbra{2j,2j}{2j,2j},
\]
with weights $p_j>0$ summing to $1$. The first has Schmidt rank $r$, the
second is a mixture of product states and is separable. A quadrature monomial
of degree two changes the photon numbers by at most two in total, while
passing from $\ket{2j,2j}$ to $\ket{2j',2j'}$ with $j\ne j'$ requires changing
each photon number by at least two, four quanta in all. The off-diagonal
coherences of $\ket{\psi_r}$, which carry its entanglement, are therefore
invisible to every first and second moment, and also to third moments. With
geometric weights $p_j\propto z^j$ the mean photon number stays below
$4z/(1-z)$ for every $r$.

The main result (Theorem~\ref{thm:main} and Corollary~\ref{cor:uniform}) states
the general form of this observation: for every degree $d$, every energy bound
$E>0$ and every $r\ge2$, there is a pure state of Schmidt rank $r$ and a
separable state whose quadrature moments of degree at most $d$ all coincide and
whose common mean photon number is at most $E$. The construction uses the
spacing $m$ with $2m>d$ in place of $2$, and Proposition~\ref{prop:sharp}
shows that this condition cannot be relaxed for these states. The obstruction
concerns uniform identification of the Schmidt number from low-order moments.
It does not prevent a particular covariance matrix from certifying a Schmidt
number larger than one, as happens for the Gaussian and photon-added states
studied in~\cite{Mishra2026}. Proposition~\ref{prop:markov} shows that at
bounded energy the obstruction necessarily involves a small Schmidt
coefficient, and Section~\ref{sec:scope} formulates the more stable question
that this suggests.

\section{Related work}\label{sec:related}
Duan, Giedke, Cirac and Zoller~\cite{Duan2000} and Simon~\cite{Simon2000}
gave separability criteria for two modes in terms of second moments; for
Gaussian states of two modes these criteria are necessary and sufficient, and
for general states they are sufficient conditions for entanglement. Shchukin
and Vogel~\cite{Shchukin2005} extended this approach to a hierarchy of
conditions on moments of all orders, which together characterize the states
with a non-positive partial transpose. Our example sits at the interface of
these works: $\ket{\psi_r}$ has a non-positive partial transpose, but its
moments of degree below $2m$ coincide with those of a separable state, so only
conditions involving moments of degree at least $2m$ can detect it.

The Schmidt number of a mixed state was introduced by Terhal and
Horodecki~\cite{Terhal2000} as the smallest $k$ such that the state is a
mixture of pure states of Schmidt rank at most $k$. Mishra
et~al.~\cite{Mishra2026} bound the fidelity of a continuous-variable state
with a pure Gaussian target, typically a two-mode squeezed vacuum, using only
the covariance matrices of the two states and the mean excitation number, and
turn this bound into Schmidt-number witnesses; they compare it with bounds
from the entanglement of formation and with linear Gaussian witnesses, on
two-mode squeezed thermal states and on photon-added displaced squeezed states.
Every criterion of this type takes the same value on $\ket{\psi_r}$ and on the
separable $\sigma_r$, and therefore certifies no entanglement for
$\ket{\psi_r}$; the result below makes this limitation uniform in $r$ at a
fixed energy. In finite dimension, Yi, Li and Zhu~\cite{Yi2026} certify
entanglement dimensionality from moments of a different kind, the moments of
the operator obtained by applying the $k$-reduction map to the state, which are
estimated by randomized measurements. Those moments are nonlinear functionals
of the state and are not determined by low-order quadrature moments, so our
obstruction does not apply to them.

\section{Results}\label{sec:results}
Let $\hat a,\hat b$ be the annihilation operators of the two modes and
$\hat x_A=(\hat a+\hat a^\dagger)/\sqrt2$,
$\hat p_A=(\hat a-\hat a^\dagger)/(i\sqrt2)$, and similarly $\hat x_B,\hat p_B$.
A quadrature moment of degree $\ell$ of a state $\rho$ is
$\Tr(\rho\,\hat W)$ for an ordered product $\hat W$ of $\ell$ operators taken
among $\hat x_A,\hat p_A,\hat x_B,\hat p_B$. Symmetrized moments, the mean
vector and the covariance matrix are linear combinations of these. All states
considered here have finite support in the Fock basis, so every moment is
finite. For integers $r\ge2$ and $m\ge1$ and weights $p_0,\dots,p_{r-1}>0$ with
$\sum_jp_j=1$, set
\[
\ket{\psi}=\sum_{j=0}^{r-1}\sqrt{p_j}\,\ket{mj,mj},
\qquad
\sigma=\sum_{j=0}^{r-1}p_j\,\ketbra{mj,mj}{mj,mj}.
\]

\begin{theorem}\label{thm:main}
The state $\ket\psi$ has Schmidt rank $r$, and $\sigma$ is separable. Every
quadrature moment of degree $\ell<2m$ takes the same value on $\ketbra\psi\psi$
and on $\sigma$. The mean total photon number of both states equals
$2m\sum_jj\,p_j$.
\end{theorem}

\begin{proof}
The vectors $\ket{mj}$ are orthonormal on each side, so the given expansion of
$\ket\psi$ is a Schmidt decomposition with $r$ nonzero coefficients, and
$\sigma$ is a convex combination of product states. Since
$\hat x_A,\hat p_A$ are linear combinations of $\hat a,\hat a^\dagger$, and
similarly on $B$, an ordered product $\hat W$ of $\ell$ quadratures is a linear
combination of ordered products of $\ell$ operators among
$\hat a,\hat a^\dagger,\hat b,\hat b^\dagger$. Such a product maps $\ket{n_A,n_B}$
to a multiple of $\ket{n_A+s_A,n_B+s_B}$ with $\lvert s_A\rvert+\lvert s_B\rvert\le\ell$.
For $j\ne j'$ the matrix element $\bra{mj,mj}\hat W\ket{mj',mj'}$ can be nonzero
only if $s_A=s_B=m(j-j')$, which requires $2m\lvert j-j'\rvert\le\ell$ and is
impossible when $\ell<2m$. Therefore
\[
\bra\psi\hat W\ket\psi=\sum_{j,j'}\sqrt{p_jp_{j'}}\,\bra{mj,mj}\hat W\ket{mj',mj'}
=\sum_jp_j\bra{mj,mj}\hat W\ket{mj,mj}=\Tr(\sigma\hat W).
\]
The total photon number $\hat a^\dagger\hat a+\hat b^\dagger\hat b$ is diagonal
in the Fock basis, so the same computation applies to it and gives
$\sum_jp_j\cdot2mj$ for both states.
\end{proof}

For $m=2$ the theorem says that the mean vectors, the covariance matrices and
all third moments of $\ket{\psi}$ and $\sigma$ coincide; the mean vectors
vanish. The energy is the mean photon number up to the vacuum contribution, and
any Hamiltonian quadratic in the quadratures also has the same expectation in
both states.

\begin{corollary}\label{cor:uniform}
Let $d\ge1$ be an integer and $E>0$. For every $r\ge2$ there are a pure state
of Schmidt rank $r$ and a separable state whose quadrature moments of degree at
most $d$ all coincide and whose common mean photon number is at most $E$.
\end{corollary}

\begin{proof}
Choose $m=\lfloor d/2\rfloor+1$, so that $2m>d$, and $z\in(0,1)$ with
$2mz/(1-z)\le E$. Take $p_j=z^j/\sum_{i<r}z^i$ in Theorem~\ref{thm:main}. The
mean $\mu_r=\sum_{j<r}j\,p_j$ of this truncated geometric law increases with
$r$, because passing from $r$ to $r+1$ adds the value $r$, which exceeds every
value already present, so $\mu_r\le\lim_{s\to\infty}\mu_s=z/(1-z)$. The mean
photon number $2m\mu_r$ is therefore at most $2mz/(1-z)\le E$.
\end{proof}

In particular, for second moments ($d=2$, $m=2$) the bound is $4z/(1-z)$, and
any positive energy budget admits pairs of this type with arbitrarily large
Schmidt rank. Any test whose input consists of quadrature moments of degree at
most $d$ and of an upper bound on the mean photon number assigns the same
verdict to the two states, and since one of them is separable, it cannot
certify any entanglement of the other.

\begin{proposition}[Sharpness]\label{prop:sharp}
For the states of Theorem~\ref{thm:main}, the moment
$\langle\hat a^m\hat b^m\rangle$, a polynomial of degree $2m$ in the
quadratures, equals
$\sum_{j=0}^{r-2}\sqrt{p_jp_{j+1}}\,\frac{(m(j+1))!}{(mj)!}>0$ for $\ket\psi$
and $0$ for $\sigma$. In particular, for spacing $m=1$ the second moment
$\langle\hat a\hat b\rangle$ already distinguishes the two states.
\end{proposition}

\begin{proof}
Since $\hat a^m\ket{n}=\sqrt{n!/(n-m)!}\,\ket{n-m}$, the operator
$\hat a^m\hat b^m$ maps $\ket{m(j+1),m(j+1)}$ to
$\frac{(m(j+1))!}{(mj)!}\ket{mj,mj}$ and maps $\ket{0,0}$ to zero. Its diagonal
elements on the vectors $\ket{mj,mj}$ vanish, which gives $0$ for $\sigma$, and
the off-diagonal elements between consecutive vectors give the stated sum for
$\ket\psi$. The operator is a polynomial of degree $2m$ in the quadratures
because $\hat a=(\hat x_A+i\hat p_A)/\sqrt2$ and $\hat b=(\hat x_B+i\hat p_B)/\sqrt2$.
\end{proof}

At bounded energy, the uniform obstruction always rests on a small Schmidt
coefficient.

\begin{proposition}\label{prop:markov}
Let $\ket\phi$ be a pure two-mode state with mean total photon number at most
$E$, and let $\lambda_1\ge\lambda_2\ge\cdots$ be the eigenvalues of its reduced
states, the squared Schmidt coefficients. Then $\lambda_r\le E/(r(r-1))$ for
every $r\ge2$.
\end{proposition}

\begin{proof}
The mean photon number of mode $A$ is $\Tr(\rho_A\hat a^\dagger\hat a)$, where
$\rho_A=\sum_i\lambda_i\ketbra{e_i}{e_i}$ with orthonormal eigenvectors $e_i$.
Write $\hat n=\hat a^\dagger\hat a$ and $c_i=\bra{e_i}\hat n\ket{e_i}\ge0$. By Ky
Fan's minimum principle, the sum of the diagonal values of $\hat n$ on any $s$
orthonormal vectors is at least the sum of its $s$ smallest eigenvalues, so
$C_s=\sum_{i\le s}c_i\ge\sum_{i\le s}(i-1)$ for every $s$. For every $S$, Abel
summation with the nonnegative differences $\lambda_s-\lambda_{s+1}$ gives
\[
\sum_{i\le S}\lambda_ic_i=\sum_{s<S}(\lambda_s-\lambda_{s+1})C_s+\lambda_SC_S
\ge\sum_{i\le S}\lambda_i(i-1),
\]
and letting $S$ grow yields $\Tr(\rho_A\hat n)=\sum_i\lambda_ic_i\ge\sum_i\lambda_i(i-1)$. The same bound holds for mode $B$, whose
reduced state has the same nonzero eigenvalues. Hence
$E\ge2\sum_{i\le r}\lambda_i(i-1)\ge2\lambda_r\sum_{i\le r}(i-1)=\lambda_r\,r(r-1)$.
\end{proof}

For the states of Corollary~\ref{cor:uniform} the smallest nonzero
eigenvalue of the reduced state is in fact $p_{r-1}\propto z^{r-1}$,
exponentially small in $r$.

\section{Verification}\label{sec:verification}
The results are proved above. The public
\href{https://github.com/jpbruneton/open_quantum_problems/blob/main/research/entries/schmidt-number-separable-covariance/checks/verify_moments.py}{Python checker}
confirms them on examples, and its
\href{https://github.com/jpbruneton/open_quantum_problems/blob/main/research/entries/schmidt-number-separable-covariance/data/verification.json}{recorded output}
is part of the same entry. It builds the quadrature operators as Fock matrices
truncated at a level that the words considered cannot reach, so that no
truncation error occurs, and evaluates every ordered word of length at most
$d$ in $\hat x_A,\hat p_A,\hat x_B,\hat p_B$ on $\ket\psi$ and on $\sigma$: for
spacing $m=2$, degree $d=3$ and $2\le r\le6$, and for $m\in\{3,4\}$,
$d=2m-1$ and $2\le r\le5$, each with geometric and with random weights. The
largest difference is below $3\times10^{-13}$, at the level of rounding. It
also checks the value of $\langle\hat a^m\hat b^m\rangle$ in
Proposition~\ref{prop:sharp} for $m\le4$, the difference of second moments for
spacing $1$, the Schmidt rank of $\ket\psi$, the equality of mean photon
numbers, and the bound $2mz/(1-z)$ for $r\le60$, and it tests
Proposition~\ref{prop:markov} on $300$ random pure states with finite Fock
support. The bound of Proposition~\ref{prop:markov} is attained by the
maximally entangled state $r^{-1/2}\sum_{i<r}\ket{i,i}$. The run is deterministic and
takes under half a minute.

\section{Scope and limitations}\label{sec:scope}
Theorem~\ref{thm:main}, Corollary~\ref{cor:uniform} and
Propositions~\ref{prop:sharp} and~\ref{prop:markov} are proved analytically;
the checker is a floating-point confirmation on finitely many cases. The
result rules out only criteria that use quadrature moments of bounded degree,
together with an energy bound, uniformly over all states. It says nothing
against the certification of a Schmidt number for a particular state or family
of states, such as the Gaussian states for which~\cite{Mishra2026} obtains
large certified values, and it does not concern nonlinear data such as purity
or the reduction moments of~\cite{Yi2026}. The separable partner $\sigma_r$ is
mixed; we have not examined whether a pure product state, or a pure state of
smaller Schmidt rank, can share the low-order moments of $\ket{\psi_r}$.

Proposition~\ref{prop:markov} shows that at bounded energy the obstruction
needs a reduced-state eigenvalue $\lambda_r$ of at most $E/(r(r-1))$. A better-posed question
is therefore the following: given an energy bound and a lower bound on the
$r$-th Schmidt coefficient, or on the weight of a state beyond Schmidt rank
$r-1$, which order of quadrature moments suffices to certify Schmidt number at
least $r$? We leave it open.

The construction is elementary, and we make no claim of priority for it. Our
search was limited to the works cited here, in which we did not find it
stated; observations of the same kind may well appear elsewhere in the
literature on moment-based entanglement criteria.

\begingroup\small
\begin{thebibliography}{9}
\setlength{\itemsep}{2pt}
\bibitem{Duan2000}
L.-M. Duan, G. Giedke, J. I. Cirac and P. Zoller,
\emph{Inseparability Criterion for Continuous Variable Systems},
Phys. Rev. Lett. \textbf{84}, 2722 (2000);
\href{https://doi.org/10.1103/PhysRevLett.84.2722}{doi:10.1103/PhysRevLett.84.2722};
\href{https://arxiv.org/abs/quant-ph/9908056}{arXiv:quant-ph/9908056}.

\bibitem{Simon2000}
R. Simon,
\emph{Peres-Horodecki Separability Criterion for Continuous Variable Systems},
Phys. Rev. Lett. \textbf{84}, 2726 (2000);
\href{https://doi.org/10.1103/PhysRevLett.84.2726}{doi:10.1103/PhysRevLett.84.2726};
\href{https://arxiv.org/abs/quant-ph/9909044}{arXiv:quant-ph/9909044}.

\bibitem{Terhal2000}
B. M. Terhal and P. Horodecki,
\emph{Schmidt number for density matrices},
Phys. Rev. A \textbf{61}, 040301 (2000);
\href{https://doi.org/10.1103/PhysRevA.61.040301}{doi:10.1103/PhysRevA.61.040301};
\href{https://arxiv.org/abs/quant-ph/9911117}{arXiv:quant-ph/9911117}.

\bibitem{Shchukin2005}
E. Shchukin and W. Vogel,
\emph{Inseparability Criteria for Continuous Bipartite Quantum States},
Phys. Rev. Lett. \textbf{95}, 230502 (2005);
\href{https://doi.org/10.1103/PhysRevLett.95.230502}{doi:10.1103/PhysRevLett.95.230502};
\href{https://arxiv.org/abs/quant-ph/0508132}{arXiv:quant-ph/0508132}.

\bibitem{Mishra2026}
I. Mishra, S. Morelli, K. Baksov\'a, P. Rembold and N. Friis,
\emph{Certification of high-dimensional entanglement in continuous-variable systems},
preprint (2026);
\href{https://arxiv.org/abs/2610.03113}{arXiv:2610.03113}.

\bibitem{Yi2026}
C. Yi, X. Li and H. Zhu,
\emph{Certifying Entanglement Dimensionality by $k$-Reduction Moments},
PRX Quantum \textbf{7}, 010356 (2026);
\href{https://doi.org/10.1103/cc1n-gmj1}{doi:10.1103/cc1n-gmj1};
\href{https://arxiv.org/abs/2501.15360}{arXiv:2501.15360}.
\end{thebibliography}
\endgroup
\end{document}
